% Generated by publication/build_sections.py; do not edit by hand.
\paragraph{Informal verdict.}
\textbf{A --- Complete and apparently correct}, with \textbf{high confidence}.

\paragraph{Lean coverage.}
Full canonical asymptotic coverage. The preserved base classifies every continuous solution for |q|<1, handles q=0 and the zero solution, proves the two exact logarithmic-phase identities, profile symmetry/positivity/endpoint zeros and the global obstruction to a phase-free ratio equivalent. Q839Completion.lean derives qPoch continuity, the denominator and reciprocal limits, continuity of phase and phaseNeg, exact positive-zero iff classifications in both sign regimes, the negative-q envelope bridge, and positive- and negative-q fixed-phase ratio limits and IsEquivalent theorems for every u0 in (0,1) and every nonzero continuous solution.

\subsection{Audited informal answer}
Q839 --- Complete solution
\begin{enumerate}
\item Faithful restatement
\end{enumerate}
Let \(q\in(-1,1)\) be real, and let \(f:\mathbb R\to\mathbb R\) be continuous and satisfy

\[
f(x)=(1-qx)f(qx)\qquad(x\in\mathbb R).
\]

The preceding oral exercise established that \(f\) has a power-series expansion valid on all of \(\mathbb R\). Q839 asks for the asymptotic behavior of \(f(x)\) as \(x\to+\infty\). There are no subparts, footnotes, or apparent typographical ambiguities attached to Q839 in the PDF. Luc Pommeret+1
Because nonzero solutions have infinitely many positive zeros, no ordinary global equivalent \(f(x)\sim g(x)\) with \(g\) eventually nonzero can exist. The correct result is an exact phase-dependent, logarithmically periodic asymptotic formula.

\begin{enumerate}
\item Statement of the result
\end{enumerate}
For \(|q|<1\), introduce the finite and infinite \(q\)-Pochhammer products

\[
(a;q)_N:=\prod_{k=0}^{N-1}(1-aq^k),
\qquad
(a;q)_\infty:=\prod_{k=0}^{\infty}(1-aq^k),
\]

and in particular

\[
(q;q)_N=\prod_{k=1}^{N}(1-q^k),\qquad (q;q)_0=1.
\]

Classification of all continuous solutions
Let \(c=f(0)\).


If \(q=0\), then \(f(x)=c\) for every \(x\).


If \(0<|q|<1\), then



\[
\boxed{\;
f(x)=c\,(qx;q)_\infty
=c\prod_{j=1}^{\infty}(1-q^j x).
\;}
\]

In particular, \(f\) extends to an entire function and

\[
\boxed{\;
f(x)
=
c\sum_{m=0}^{\infty}
\frac{(-1)^m q^{m(m+1)/2}}{(q;q)_m}\,x^m.
\;}
\]

Now suppose \(q\neq0\), and put

\[
a:=\log\frac1{|q|}>0,
\qquad
\mathcal E_q(x)
:=
x^{-1/2}
\exp\!\left(\frac{(\log x)^2}{2a}\right).
\]

This is the universal growth envelope.

Case \(0<q<1\)
For \(x>1\), write

\[
\frac{\log x}{\log(1/q)}=n+u,
\qquad
n\in\mathbb N,\quad 0\le u<1.
\]

Equivalently,

\[
n=\left\lfloor\frac{\log x}{\log(1/q)}\right\rfloor,
\qquad
u=\left\{\frac{\log x}{\log(1/q)}\right\}.
\]

Define

\[
\Phi_q(u)
:=
q^{(u^2-u)/2}
(q^u;q)_\infty(q^{1-u};q)_\infty,
\qquad 0\le u\le1.
\]

Then the following identity is exact for every \(x>1\):

\[
\boxed{\;
f(x)
=
c(-1)^n\mathcal E_q(x)
\frac{\Phi_q(u)}{(x^{-1};q)_\infty}.
\;}
\tag{2.1}
\]

Moreover,

\[
\Phi_q(0)=\Phi_q(1)=0,\qquad
\Phi_q(u)>0\quad(0<u<1),
\]

and \(\Phi_q(1-u)=\Phi_q(u)\).

Case \(-1<q<0\)
Set

\[
r:=|q|=-q,\qquad s:=r^2=q^2.
\]

For \(x>1\), write

\[
\frac{\log x}{\log(1/s)}=n+u,
\qquad
n\in\mathbb N,\quad 0\le u<1.
\]

Thus

\[
n=\left\lfloor\frac{\log x}{2\log(1/r)}\right\rfloor,
\qquad
u=\left\{\frac{\log x}{2\log(1/r)}\right\}.
\]

Define

\[
\begin{aligned}
\Psi_q(u)
:={}&
s^{u^2-u/2}
(s^u;s)_\infty(s^{1-u};s)_\infty\\
&\times
(-s^{u+1/2};s)_\infty
(-s^{1/2-u};s)_\infty,
\qquad 0\le u\le1.
\end{aligned}
\]

Then, again, the following identity is exact for every \(x>1\):

\[
\boxed{\;
f(x)
=
c(-1)^n\mathcal E_q(x)
\frac{\Psi_q(u)}{(x^{-1};q)_\infty}.
\;}
\tag{2.2}
\]

Here

\[
\Psi_q(0)=\Psi_q(1)=0,\qquad
\Psi_q(u)>0\quad(0<u<1),
\]

and \(\Psi_q(1-u)=\Psi_q(u)\).

Complete asymptotic expansion
For every fixed \(q\in(-1,1)\setminus\{0\}\) and \(x>1\),

\[
\boxed{\;
\frac1{(x^{-1};q)_\infty}
=
\sum_{m=0}^{\infty}\frac{x^{-m}}{(q;q)_m}.
\;}
\tag{2.3}
\]

Consequently, in either sign case, if \(A_q=\Phi_q\) or \(\Psi_q\) and \(n,u\) are defined as above, then

\[
\boxed{\;
f(x)
=
c(-1)^n\mathcal E_q(x)A_q(u)
\sum_{m=0}^{\infty}\frac{x^{-m}}{(q;q)_m}.
\;}
\tag{2.4}
\]

Thus, for every integer \(N\ge1\),

\[
\boxed{\;
f(x)
=
c(-1)^n\mathcal E_q(x)A_q(u)
\left(
\sum_{m=0}^{N-1}\frac{x^{-m}}{(q;q)_m}
+O_{q,N}(x^{-N})
\right).
\;}
\tag{2.5}
\]

The error estimate is uniform over all logarithmic phases \(u\in[0,1)\). In particular,

\[
f(x)
=
c(-1)^n\mathcal E_q(x)A_q(u)
\left(
1+\frac{1}{(1-q)x}
+\frac{1}{(1-q)(1-q^2)x^2}
+O_q(x^{-3})
\right).
\]

All constants here are for fixed \(q\); uniformity as \(q\to0\) or \(|q|\to1\) is not asserted.

\begin{enumerate}
\item Classification of the solutions
\end{enumerate}
Iterating the functional equation gives, for every \(N\ge1\),

\[
f(x)
=
\prod_{j=1}^{N}(1-q^j x)\,f(q^N x).
\tag{3.1}
\]

Since \(|q|<1\), we have \(q^N x\to0\). By continuity at \(0\),

\[
f(q^N x)\longrightarrow f(0)=c.
\]

For fixed \(x\),

\[
\sum_{j=1}^{\infty}|q^j x|<\infty,
\]

so the product

\[
\prod_{j=1}^{\infty}(1-q^j x)
\]

converges. Passing to the limit in (3.1) gives

\[
f(x)=c\prod_{j=1}^{\infty}(1-q^j x)=c(qx;q)_\infty.
\]

Conversely, this product satisfies

\[
\begin{aligned}
(1-qx)(q^2x;q)_\infty
&=(1-qx)\prod_{j=2}^{\infty}(1-q^j x)\\
&=\prod_{j=1}^{\infty}(1-q^j x),
\end{aligned}
\]

so it is indeed a solution.
The convergence is uniform on every compact subset of \(\mathbb C\), since for \(|x|\le R\),

\[
\sum_{j=1}^{\infty}|q^j x|
\le R\sum_{j=1}^{\infty}|q|^j<\infty.
\]

Thus the product defines an entire function.
To obtain the coefficients, write

\[
f(x)=\sum_{m=0}^{\infty}a_mx^m.
\]

Substitution into

\[
f(x)=(1-qx)f(qx)
\]

gives, for \(m\ge1\),

\[
a_m=q^ma_m-q^ma_{m-1},
\]

hence

\[
(1-q^m)a_m=-q^ma_{m-1}.
\]

Starting from \(a_0=c\),

\[
a_m
=
c\,\frac{(-1)^m q^{1+2+\cdots+m}}
{\prod_{j=1}^{m}(1-q^j)}
=
c\,\frac{(-1)^m q^{m(m+1)/2}}{(q;q)_m}.
\]

Furthermore,

\[
\left\lvert \frac{a_{m+1}}{a_m}\right\rvert 
=
\frac{|q|^{m+1}}{|1-q^{m+1}|}
\longrightarrow0,
\]

so the radius of convergence is infinite.

\begin{enumerate}
\item The reciprocal-product expansion
\end{enumerate}
Let

\[
G(z):=\frac1{(z;q)_\infty},
\qquad |z|<1.
\]

Because

\[
(z;q)_\infty=(1-z)(qz;q)_\infty,
\]

we have

\[
(1-z)G(z)=G(qz).
\]

Write

\[
G(z)=\sum_{m=0}^{\infty}b_mz^m.
\]

Comparison of coefficients gives \(b_0=1\), and for \(m\ge1\),

\[
b_m-b_{m-1}=q^mb_m.
\]

Therefore

\[
(1-q^m)b_m=b_{m-1},
\]

so

\[
b_m=\frac1{(q;q)_m}.
\]

This proves

\[
\frac1{(z;q)_\infty}
=
\sum_{m=0}^{\infty}\frac{z^m}{(q;q)_m},
\qquad |z|<1.
\]

Taking \(z=x^{-1}\) proves (2.3).
For completeness, \((q;q)_m\) converges as \(m\to\infty\) to the nonzero number \((q;q)_\infty\). Hence there is a constant \(C_q\) such that

\[
\frac1{|(q;q)_m|}\le C_q
\qquad(m\ge0).
\]

For \(x\ge2\),

\[
\left\lvert 
\sum_{m=N}^{\infty}\frac{x^{-m}}{(q;q)_m}
\right\rvert 
\le
C_q\frac{x^{-N}}{1-x^{-1}}
\le 2C_qx^{-N},
\]

which proves the uniform remainder estimate in (2.5).

\begin{enumerate}
\item Derivation for \(0<q<1\)
\end{enumerate}
Put

\[
a=\log\frac1q,
\qquad
\frac{\log x}{a}=n+u,
\quad 0\le u<1.
\]

Thus

\[
x=q^{-(n+u)}.
\]

Split the product after its first \(n\) factors:

\[
(qx;q)_\infty
=
(qx;q)_n\,(q^{n+1}x;q)_\infty.
\]

Since \(q^{n+1}x=q^{1-u}\),

\[
(q^{n+1}x;q)_\infty=(q^{1-u};q)_\infty.
\]

For the finite product,

\[
\begin{aligned}
(qx;q)_n
&=\prod_{j=1}^{n}(1-q^jx)\\
&=(-1)^nx^nq^{n(n+1)/2}
\prod_{j=1}^{n}\left(1-\frac{q^{-j}}x\right).
\end{aligned}
\]

Because \(x=q^{-(n+u)}\),

\[
\frac{q^{-j}}x=q^{n+u-j}.
\]

Reversing the order of the factors,

\[
\prod_{j=1}^{n}\left(1-\frac{q^{-j}}x\right)
=
\prod_{k=0}^{n-1}(1-q^{u+k})
=
(q^u;q)_n.
\]

Now

\[
(q^u;q)_n
=
\frac{(q^u;q)_\infty}{(q^{n+u};q)_\infty}
=
\frac{(q^u;q)_\infty}{(x^{-1};q)_\infty}.
\]

Therefore

\[
(qx;q)_\infty
=
(-1)^nx^nq^{n(n+1)/2}
\frac{(q^u;q)_\infty(q^{1-u};q)_\infty}
{(x^{-1};q)_\infty}.
\tag{5.1}
\]

It remains to rewrite the elementary prefactor. Since \(\log x=a(n+u)\),

\[
\begin{aligned}
\log\!\left(x^nq^{n(n+1)/2}\right)
&=an(n+u)-\frac a2n(n+1)\\
&=\frac{(\log x)^2}{2a}
-\frac12\log x
+\frac a2u(1-u).
\end{aligned}
\]

Because

\[
e^{a u(1-u)/2}=q^{(u^2-u)/2},
\]

we get

\[
x^nq^{n(n+1)/2}
=
\mathcal E_q(x)\,q^{(u^2-u)/2}.
\]

Substituting into (5.1) proves (2.1).

\begin{enumerate}
\item Derivation for \(-1<q<0\)
\end{enumerate}
Write

\[
q=-r,\qquad 0<r<1,\qquad s=r^2.
\]

Separating odd and even powers of \(q\) gives

\[
\begin{aligned}
(qx;q)_\infty
&=(-rx;-r)_\infty\\
&=(sx;s)_\infty(-rx;s)_\infty.
\end{aligned}
\tag{6.1}
\]

Let

\[
\frac{\log x}{\log(1/s)}=n+u,
\qquad 0\le u<1.
\]

Thus \(x=s^{-(n+u)}\).
Applying the positive-base calculation to the first factor gives

\[
(sx;s)_\infty
=
(-1)^nx^ns^{n(n+1)/2}
\frac{(s^u;s)_\infty(s^{1-u};s)_\infty}
{(x^{-1};s)_\infty}.
\tag{6.2}
\]

For the second factor,

\[
(-rx;s)_\infty
=
\prod_{k=0}^{n-1}(1+rxs^k)
\,(-rxs^n;s)_\infty.
\]

Since \(rxs^n=s^{1/2-u}\),

\[
(-rxs^n;s)_\infty=(-s^{1/2-u};s)_\infty.
\]

Also,

\[
\begin{aligned}
\prod_{k=0}^{n-1}(1+rxs^k)
&=
r^nx^ns^{n(n-1)/2}
\prod_{k=0}^{n-1}
\left(1+\frac1{rxs^k}\right)\\
&=
r^nx^ns^{n(n-1)/2}
(-s^{u+1/2};s)_n.
\end{aligned}
\]

Indeed,

\[
\frac1{rxs^k}
=
s^{n+u-k-1/2},
\]

whose exponents, in reverse order, are

\[
u+\frac12,\ u+\frac32,\ldots,u+n-\frac12.
\]

Moreover,

\[
(-s^{u+1/2};s)_n
=
\frac{(-s^{u+1/2};s)_\infty}
{(-s^{n+u+1/2};s)_\infty}
=
\frac{(-s^{u+1/2};s)_\infty}
{(-r/x;s)_\infty}.
\]

Hence

\[
(-rx;s)_\infty
=
r^nx^ns^{n(n-1)/2}
\frac{
(-s^{u+1/2};s)_\infty
(-s^{1/2-u};s)_\infty
}{
(-r/x;s)_\infty
}.
\tag{6.3}
\]

Multiplying (6.2) and (6.3), the denominator is

\[
(x^{-1};s)_\infty(-r/x;s)_\infty
=
(x^{-1};-r)_\infty
=
(x^{-1};q)_\infty.
\]

The elementary prefactor is

\[
(-1)^nx^{2n}r^ns^{n^2}.
\]

Because \(r=s^{1/2}\),

\[
x^{2n}r^ns^{n^2}
=
x^{2n}s^{n^2+n/2}.
\]

Let \(A=\log(1/s)=2a\). Since \(\log x=A(n+u)\),

\[
\begin{aligned}
\log\!\left(x^{2n}s^{n^2+n/2}\right)
&=
A\left(n^2+2nu-\frac n2\right)\\
&=
\frac{(\log x)^2}{A}
-\frac12\log x
+A\left(\frac u2-u^2\right).
\end{aligned}
\]

Now \(A=2a\), and

\[
e^{A(u/2-u^2)}
=
s^{u^2-u/2}.
\]

Therefore

\[
x^{2n}r^ns^{n^2}
=
\mathcal E_q(x)s^{u^2-u/2}.
\]

This proves (2.2).

\begin{enumerate}
\item Properties of the phase profiles
\end{enumerate}
All factors defining \(\Phi_q\) and \(\Psi_q\) converge uniformly for \(u\in[0,1]\), after separating finitely many initial factors when necessary. Thus both profiles are continuous.
For \(0<u<1\), every factor in \(\Phi_q(u)\) is positive, so \(\Phi_q(u)>0\). At \(u=0\), the factor \((q^0;q)_\infty=(1;q)_\infty\) vanishes; at \(u=1\), the other Pochhammer factor vanishes. Symmetry is immediate.
For \(\Psi_q\), every factor is also positive in the open interval. Its symmetry needs one short computation. Put

\[
B(u):=(-s^{u+1/2};s)_\infty(-s^{1/2-u};s)_\infty.
\]

Using

\[
(-z;s)_\infty=(1+z)(-sz;s)_\infty,
\]

one obtains

\[
\frac{B(1-u)}{B(u)}
=
\frac{1+s^{u-1/2}}{1+s^{1/2-u}}
=
s^{u-1/2}.
\]

On the other hand,

\[
\frac{s^{(1-u)^2-(1-u)/2}}{s^{u^2-u/2}}
=
s^{1/2-u}.
\]

The two factors cancel, proving

\[
\Psi_q(1-u)=\Psi_q(u).
\]

For \(q>0\), the maximum of \(\Phi_q\) is explicit. Indeed, writing \(q=e^{-a}\), one has on \(0<u<1\)

\[
\begin{aligned}
\log\Phi_q(u)
={}&
\frac a2u(1-u)\\
&+\sum_{k=0}^{\infty}
\log(1-e^{-a(k+u)})
+\sum_{k=0}^{\infty}
\log(1-e^{-a(k+1-u)}).
\end{aligned}
\]

Each logarithmic summand has negative second derivative, and the first term has second derivative \(-a\). Thus \(\log\Phi_q\) is strictly concave. By symmetry its unique maximum is at \(u=1/2\), so

\[
\boxed{\;
\max_{0\le u\le1}\Phi_q(u)
=
q^{-1/8}(q^{1/2};q)_\infty^2.
\;}
\]

For negative \(q\), the exact continuous profile \(\Psi_q\) already gives the extremal normalized amplitude; no unimodality assertion is needed.

\begin{enumerate}
\item Zeros, signs, oscillations, and growth
\end{enumerate}
Assume \(c\neq0\).
From the product representation, all zeros are

\[
x=q^{-j},\qquad j=1,2,\ldots,
\]

and they are simple.
Thus the positive zeros are

\[
\begin{cases}
q^{-j},&0<q<1,\\[1mm]
|q|^{-2j},&-1<q<0,
\end{cases}
\qquad j\ge1.
\]

For \(q>0\),

\[
\operatorname{sgn}f(x)
=
\operatorname{sgn}(c)(-1)^n
\]

on

\[
q^{-n}<x<q^{-(n+1)}.
\]

For \(q<0\),

\[
\operatorname{sgn}f(x)
=
\operatorname{sgn}(c)(-1)^n
\]

on

\[
|q|^{-2n}<x<|q|^{-2(n+1)}.
\]

Let

\[
M_q:=
\begin{cases}
\displaystyle\max_{0\le u\le1}\Phi_q(u),&q>0,\\[2mm]
\displaystyle\max_{0\le u\le1}\Psi_q(u),&q<0.
\end{cases}
\]

Since \((x^{-1};q)_\infty\to1\), the complete cluster set of the normalized function is

\[
\boxed{\;
\operatorname{Clust}_{x\to\infty}
\frac{f(x)}{c\,\mathcal E_q(x)}
=
[-M_q,M_q].
\;}
\]

Indeed, every sequence has a subsequence along which the logarithmic phase \(u\) converges and the parity of \(n\) is fixed. Conversely, every value in the interval is obtained because the phase profile is continuous, vanishes at the endpoints, and attains \(M_q\).
Consequently,

\[
\boxed{\;
\limsup_{x\to\infty}
\frac{|f(x)|}{\mathcal E_q(x)}
=
|c|M_q,
\qquad
\liminf_{x\to\infty}
\frac{|f(x)|}{\mathcal E_q(x)}
=
0.
\;}
\]

The second equality follows already from the positive zeros. Moreover,

\[
\boxed{\;
\limsup_{x\to\infty}f(x)=+\infty,
\qquad
\liminf_{x\to\infty}f(x)=-\infty.
\;}
\]

Finally,

\[
\log \mathcal E_q(x)
=
\frac{(\log x)^2}{2\log(1/|q|)}
-\frac12\log x.
\]

Hence

\[
\boxed{\;
\limsup_{x\to\infty}
\frac{\log(1+|f(x)|)}{(\log x)^2}
=
\frac{1}{2\log(1/|q|)},
\qquad
\liminf_{x\to\infty}
\frac{\log(1+|f(x)|)}{(\log x)^2}
=
0.
\;}
\]

Thus the nonzero solution has superpolynomial peak growth but returns exactly to zero along a geometric sequence.
For example, at any fixed phase \(u_0\in(0,1)\),

\[
f\!\left(q^{-(n+u_0)}\right)
\sim
c(-1)^n\Phi_q(u_0)
\mathcal E_q\!\left(q^{-(n+u_0)}\right)
\]

when \(q>0\), while for \(q<0\),

\[
f\!\left((q^2)^{-(n+u_0)}\right)
\sim
c(-1)^n\Psi_q(u_0)
\mathcal E_q\!\left((q^2)^{-(n+u_0)}\right).
\]


\begin{enumerate}
\item Boundary and hypothesis checks
\end{enumerate}
The case \(q=0\)
The equation becomes

\[
f(x)=f(0),
\]

so every solution is constant. The logarithmic-envelope formulas are not meant to include this case.
The zero solution
If \(f(0)=0\), the classification gives \(f\equiv0\).
Two elementary examples
For \(q=\tfrac12\),

\[
f(x)=c\prod_{j=1}^{\infty}(1-2^{-j}x),
\]

with positive zeros \(2,4,8,\ldots\), logarithmic phase \(\{\log_2x\}\), and envelope

\[
x^{-1/2}\exp\!\left(\frac{(\log x)^2}{2\log2}\right).
\]

For \(q=-\tfrac12\),

\[
f(x)=c\prod_{j=1}^{\infty}(1-(-2)^{-j}x),
\]

with positive zeros \(4,16,64,\ldots\), phase \(\{\log_4x\}\), and the same envelope.
Why continuity matters
Only continuity at \(0\) is actually needed. Without it, the solution space is much larger.
For \(q\neq0\), choose \(x_0\neq0\) whose multiplicative orbit

\[
\mathcal O=\{q^kx_0:k\in\mathbb Z\}
\]

does not contain \(q^{-1}\). Choose \(v_0=1\), and define \(v_k\) recursively by

\[
v_k=(1-q^{k+1}x_0)v_{k+1}.
\]

Set

\[
g(q^kx_0)=v_k,\qquad g(x)=0\quad(x\notin\mathcal O).
\]

The orbit and its complement are both invariant under multiplication by \(q\), so \(g\) satisfies the functional equation. But as \(k\to+\infty\),

\[
q^kx_0\to0,
\qquad
v_k\longrightarrow
\frac1{(qx_0;q)_\infty}\neq0,
\]

while one may set \(g(0)=0\). Thus \(g\) is discontinuous at \(0\).
Why the strict condition \(|q|<1\) changes the problem
At \(q=1\),

\[
f(x)=(1-x)f(x),
\]

so \(xf(x)=0\); continuity forces \(f\equiv0\).
At \(q=-1\),

\[
f(x)=(1+x)f(-x),\qquad
f(-x)=(1-x)f(x),
\]

and hence

\[
f(x)=(1-x^2)f(x).
\]

Thus \(x^2f(x)=0\), and continuity again forces \(f\equiv0\).

Final conclusion
Every continuous solution is a scalar multiple of the Euler product

\[
\prod_{j\ge1}(1-q^jx).
\]

For \(q\neq0\), its natural magnitude is

\[
x^{-1/2}
\exp\!\left(
\frac{(\log x)^2}{2\log(1/|q|)}
\right),
\]

but the normalized function does not converge: it has an explicit logarithmically periodic profile, changes sign after each positive zero, and admits the exact convergent inverse-power expansion (2.4). This completely determines its behavior as \(x\to+\infty\).
